\subsection{The actual algebraic raw theta sections}
\label{sec:actual-raw-theta}

The normalized finite theta sum defines an algebraic section on the
original moduli space.  Its construction works over every coefficient
algebra, including nonreduced algebras.  In the actual positive weight
range it is a cusp section, with exactly the value used by the integral
denominator theorem.  Thus that theorem applies to constructed sections
in its original target space.

Keep the actual host, datum, primitive line, coefficient field $K=c.F$
and an arbitrary original $\mathrm{IsNab}$ rule.  Put $r=p+q$ and
write $\mathfrak b=\mathrm{Mb}$, so
$\mathfrak b\bar{\mathfrak b}=(M)$ and
$\mathfrak b+\bar{\mathfrak b}=\mathcal O_E$.
Let $\Lambda\subset E^n$ be an original self-dual lattice, standard
locally at $M$.  Let $N$ be positive and odd, with
$M\mid N$ and $N\mathcal O_E^n\subset\Lambda\cap\mathcal O_E^n$.
For the original row weight $a$, let $\mathcal H_a$ be the
$K$-space of harmonic row-highest polynomials in the prescribed
variables, and retain the original $\mathrm{Dom}(a,n)$ hypothesis.
Set
\[
 I=[\Lambda\cap\mathcal O_E^n:N\mathcal O_E^n]^r,
 \qquad J=[\Lambda:N\mathcal O_E^n]^r,
 \qquad C=\#(\mathrm{LevQ}^{,n})=M^{[E:\mathbb Q]rn/2}.
\]
These two indices are different and are not interchanged below.

\begin{theorem}[Exact raw-theta existence and integral image]
\label{thm:actual-raw-theta-existence}
For every function $w:\mathrm{LevQ}^{,n}\to K$ there is a
$K$-linear map
\[
 \mathcal T_{\Lambda,N,w}:\mathcal H_a
                    \longrightarrow\operatorname{sec}(a,n)
\]
with the following property.  For every $K$-algebra $S$, every
original framed point $y/S$, and every $P\in\mathcal H_a$, there
is a unique value $v_y\in S$ satisfying the original
$\mathrm{thLaValAt}(w,\Lambda,N,P,y,v_y)$ contract.
The values commute with every algebra map and every original
$\mathrm{FrMap}$.  On test rings they are exactly the values of
$\mathcal T_{\Lambda,N,w}(P)$; on every field they also agree with
the original evaluation $\mathrm{evR}$.

Under the actual $\mathrm{CovGmin.Hyp}$ conditions the image lies in
\[
       Y.V(\mathrm{wtZ}(a,n))
           =\operatorname{cuspW}(\mathrm{wtZ}(a,n)).
\]
In particular take
$w=\widehat{\mathbf1}_{x_0}$ for any original
$x_0\in\mathrm{LevQc}^{,n}$.  Under the unchanged hypotheses of
the actual denominator theorem, for every
$P\in\operatorname{harmL}(\Lambda,a)$ and every original
determinant-ideal scalar $s$,
\[
 Zs\,\mathcal T_{\Lambda,N,\widehat{\mathbf1}_{x_0}}(P)
                 \in\operatorname{latW},
 \qquad
 Z=M^{[E:\mathbb Q]rn/2+d_a}|\operatorname{disc}E|^{d_a},
 \qquad d_a=\mathrm{DF}(D_w,a).
\]
The integer $Z$ and its previously proved size bound are unchanged.
The coefficient field, primitive line, connection, separate
eigenlattices and finite level labels are the original ones.
\end{theorem}

\begin{proof}
\emph{The finite scheme on which the value is constructed.}
Over a $K$-algebra, all positive integers are units.  Hence the
$N^2$-torsion of the original abelian scheme is finite \'etale.
Inside its $n$th power take the subgroup $Z_\Lambda$ cut out by
the original annihilator equations $\mathrm{InZh}$.
Those equations involve finitely many generators of a projective
$\mathcal O_E$-module; the definition using all annihilator rows
cuts out the same subgroup.  Its geometric fibres, in the given
Tate marking, are
\[
 Z_\Lambda=(N^{-1}\Lambda\otimes_{\mathcal O_E}L)/L^n,
                                                        \tag{RT1}
\]
where $L$ is the original period lattice of rank $r$ over
$\mathcal O_E$.  To check the formula, a class represented by
$v/N^2$ satisfies the row equations precisely when every row of
the dual of $\Lambda$ sends $v/N$ into $L$.  Double duality of
projective lattices gives $v/N\in\Lambda\otimes L$.
Self-duality gives $\Lambda\subset N^{-1}\mathcal O_E^n$, so
the displayed quotient is indeed killed by $N^2$.
Localization at every prime, followed by the determinant-index
formula, gives its order $J$.  Thus $Z_\Lambda$ is finite \'etale
of constant rank $J$ over the original base.

For each $q\in\mathrm{LevQ}^{,n}$ choose, after a finite
\'etale surjective base change, roots $t_{q,b}$ with
$[N]t_{q,b}=\eta(q_b)$ that are killed by a power of
$\mathfrak b$.  Such roots exist: project an arbitrary
$N^2$-torsion root to its $\mathfrak b$-primary factor by the
Chinese remainder decomposition of the $N^2$-torsion module.
The complementary factor includes both the
$\bar{\mathfrak b}$-primary part and any primes dividing $N$
but not $M$.  The prescribed $\eta(q_b)$ is
$\mathfrak b$-primary, so this projection preserves its image
under $[N]$.
Since $[N]\eta(q_b)=0$, these roots are killed by $N^2$.
The same finite cover splits $Z_\Lambda$ into $J$ labelled
sections.  The canonical odd theta lifts and normalized pullback
of the original primitive line give, for every $q,z,b$, the
unique section
\[
          F_{q,z,b}\in H^0(A_S,\mathscr L^{N^2})
\]
specified by $\mathrm{IsUUOver}(N,y.s,z_b,t_{q,b})$.
They use the original order of the two translations and the
original pullback normalization.  Existence follows from the
cubical pullback identity and the theta-group action.  One can
check descent of the odd lifts directly.  Put $m=N^2$.
For any local lift $U$, write $a=U^m$ and
$b=U\iota(U)$, where $\iota$ is induced by inversion.
These are central units and $a^2=b^m$.  Multiplying $U$ by
$c=ab^{-(m+1)/2}$ makes its $m$th power and its product with
its image under inversion both one.  Two such lifts differ by a scalar
with $c^2=c^m=1$, hence by one.  Uniqueness descends the
normalized lift.  The original normalized pullback and ordered
action consequently give the unique section, as in the actual
integral-theta construction.\cite{raw-theta-source}

For $d=\deg P$ form the finite expression
\begin{equation}\label{eq:actual-raw-theta-sum}
 V_y(P)=\frac{1}{ICN^d}
       \sum_{q\in\mathrm{LevQ}^{,n}}w(q)
       \sum_{z\in Z_\Lambda}
          H_y\bigl(P;(F_{q,z,b})_{b=1}^n\bigr).
\end{equation}
Here $H_y$ is literally the original product of column connection
values, summed against the coefficients of $P$.  In particular
this is the expression with $\beta=1$ in
$\mathrm{ShortValAt}$, which defines $\mathrm{thLaValAt}$.

\emph{Independence of the finite choices.}
The product Weil pairing is one on
$Z_\Lambda\times Z_\Lambda$.  Indeed, for the representatives
in \textup{(RT1)}, the product pairing is the exponential of the
integral trace pairing on $\Lambda\otimes L$; integrality follows
from the original self-duality and adjoint identity.
Also the Weil pairing is one between two
$\mathfrak b$-primary points.  Choose
$b+\bar c=1$ with $b,c\in\mathfrak b^e$ annihilating those
points, and apply the adjoint identity in the two factors.
These statements hold over nonreduced coefficient algebras too:
the relevant values lie in the finite \'etale group
$\mu_{N^2}$, and equality of its sections is detected on
geometric fibres.  Equivalently, if $u^{N^2}=1$ and $u-1$ is
nilpotent, then
$1+u+\cdots+u^{N^2-1}$ is a unit, so $u=1$.

Changing $t_q$ to another allowed choice changes it by an
$N$-torsion, $\mathfrak b$-primary vector $u$.
The condition that $\Lambda$ is standard at $M$ puts $u$ in
$Z_\Lambda$: clear the lattice comparison by an integer prime
to $M$, and invert that integer on the $\mathfrak b$-primary
group in each annihilator equation.
The translation $z\mapsto z+u$ permutes the full group.
An arbitrary original witness also enumerates that full group:
its $J$ sections are geometrically distinct in every fibre of
the finite \'etale scheme of rank $J$.  Two enumerations
therefore differ by a locally constant permutation; their
finite sums can be compared on the corresponding clopen
pieces of a common cover.
The canonical odd theta lifts compare the column sections by
scalars $r_b$ satisfying
\[
       r_b^2=e_{N^2}(u_b,z_b),\qquad r_b^{N^2}=1.
\]
There is no additional factor from $t_q$, by half-primary
isotropy.  Since $\prod_b e_{N^2}(u_b,z_b)=1$ and $N$ is
odd, $\prod_b r_b=1$.  Multilinearity in the column sections
therefore identifies each reindexed product value exactly.
This proves independence of the roots and enumeration before
any harmonic value is taken; in particular it does not assume
that the individual column commutators are one.

The actual connection construction proves that harmonic product
values commute with arbitrary base change and are independent
of the allowed canonical formal chart.  The remaining changes of
gauge contribute the paired contraction operators, which kill
$P$.  Thus the two pullbacks of
\eqref{eq:actual-raw-theta-sum} to the tensor square of the finite
faithfully flat cover agree.  The faithfully flat equalizer
\[
                   S\longrightarrow S'
                         \rightrightarrows S'\otimes_S S'
\]
gives a unique $v_y\in S$.  The same comparison on a common
cover proves the original value contract for every permitted
witness and proves its uniqueness.  Pulling the construction
along any algebra map proves its stated naturality.
No reduction of $S$ or enlargement of $K$ has been made.

\emph{The weight and the algebraic section.}
For a triangular frame change with diagonal units
$(t_0,t_{\tau i})$, the theta generator contributes $t_0^n$.
The derivative directions transform by that same triangular
matrix, and the row-highest identity for $P$ contributes
$\prod_{\tau,i}t_{\tau i}^{a(\tau,i)}$.  Consequently
\[
 V_{y'}(P)=t_0^n\prod_{\tau,i}t_{\tau i}^{a(\tau,i)}V_y(P).
                                                        \tag{RT2}
\]
This is exactly $\mathrm{wchar}(\mathrm{ex}(D_w,a,n),t)$,
including the line coordinate of exponent $n$.
The naturality just proved and \textup{(RT2)} are precisely the
two original clauses of $\operatorname{sec}(a,n)$.
Equivalently, on the actual frame torsor the finite expression is
regular and descends with this character to the flag line.
All its ingredients are algebraic finite torsion, cubical
comparisons and finite formal jets.

Every nonzero polynomial in $\mathcal H_a$ has the same total
degree $d_a$.  Hence one may use $N^{-d_a}$ throughout the
linear formula; the zero polynomial contributes zero regardless
of its convention for total degree.  This proves $K$-linearity.
The original $\mathrm{AllTest}$ presentation of a field point,
together with naturality, identifies $v_y$ with $\mathrm{evR}$.
The cusp verification below puts these sections in the original
$Y.V(\mathrm{wtZ}(a,n))$ under $\mathrm{CovGmin.Hyp}$.
The actual denominator theorem then gives the last assertion with
its same $Z$ and every original valuation test.
\end{proof}

\paragraph{The classical reading with both indices retained.}
For the original even $n$ (as required by $\mathrm{Dom}(a,n)$),
at a complex framed point choose its already constructed
analytic fibre, original symplectic marking and letter dictionary.
Let $v_0$ be the norm-one finite vacuum for
$\mathcal O_E^n\otimes L$, and write the exact product of the
chosen primitive theta generator as
$\prod_b\theta(z_b)=\kappa\,\operatorname{thA}(v_0)(z)$.
Here $\kappa\ne0$ is the actual comparison scalar; it is not
reset to one.  For the original half-level roots $Y_q$ set
\[
 g\phi=\sum_q\sigma(w(q))\,\operatorname{heis}(Y_q)
             \operatorname{proj}_{\Lambda\otimes L,w_0}(v_0).
\]
Then the constructed section satisfies the exact formula
\begin{equation}\label{eq:actual-raw-theta-classical}
 \operatorname{ev}_y\mathcal T_{\Lambda,N,w}(P)
   =\sigma\!\left(\frac{J}{IC}\right)\kappa\,
             \operatorname{thP}(g\phi,P^\sigma)(x).
\end{equation}
This is the original $\beta=1$ complex reading.
For the convention with $\beta_H$ the right side is multiplied
by $\sigma(\beta_H)^{-1}$; multiplying that section by the
original $\beta_H$ recovers the literal raw value.

To verify \eqref{eq:actual-raw-theta-classical}, choose
$u_1,\ldots,u_J\in2(\Lambda\otimes L)$ representing
$(\Lambda\otimes L)/(NL^n)$; multiplication by two permutes
this quotient, whose exponent divides the odd number $N^2$.
The corresponding points of \textup{(RT1)} are $u_i/N$,
and $Nu_i\in2L^n$.  The original characteristic is one
on these even representatives.  Averaging the genuine
character-twisted lattice action therefore gives
\[
 \sum_{i=1}^J\operatorname{heis}(u_i)v_0
       =J\operatorname{proj}_{\Lambda\otimes L,w_0}v_0.
\]
Indeed the average is invariant for that lattice action, and
vectors in its intersection with $L^n$ act trivially on the
vacuum; repeated representatives contribute exactly their
common multiplicity.  Thus the average is the original
orthogonal projection, including its normalization.
Choose the original half-level representatives so that the
torsion root is $Y_q/N$ and $NY_q\in2L^n$.
The canonical ordered translation sections then read as
$\operatorname{heis}(Y_q)\operatorname{heis}(u_i)v_0$
evaluated at $Nz$, with $z$ the analytic variable.
The harmonic connection value is the original
polynomial differential operator on that product.  Its
homogeneous scaling contributes exactly $N^d$, cancelling the
$N^{-d}$ in \eqref{eq:actual-raw-theta-sum}; the source uses
ordinary derivatives, so no divided-power factorial is inserted.
The theta series and each fixed finite derivative converge
normally on compact positive-period domains, by the Gaussian
majorant.  This justifies all differentiations and finite sums.
The established relative rigidified line comparison supplies
joint holomorphy; a fibrewise metric gauge is not used as a
holomorphic trivialization.\cite{raw-theta-source}

Along each original complex component, the scalar in
\eqref{eq:actual-raw-theta-classical} is constant when the prescribed
normalized theta frames are used.  Indeed, writing $\rho(x)$ for the
single-column scalar in the Member construction, that construction gives
$\rho(x)^2=C_\xi\det h(x)$.  Since $n$ is even,
$\rho(x)^n=C_\xi^{n/2}\det h(x)^{n/2}$, and the determinant is
already the one in $\operatorname{thA}$.
The fixed finite vacuum contributes only its fixed scalar.
Thus the actual $\kappa$, including its characteristic and marking
dependence, is retained; no varying holomorphic multiplier is dropped.
The marked rational representatives $Y_q$ and the finite projection
are fixed on that component.  Consequently its reading is a finite
linear combination of the original theta series with rational
characteristics and harmonic polynomial $P$.

\paragraph{The original cusp target.}
Suppose first that $p,q>0$ and that every first-$p$ row degree is
positive.  This is exactly the extra condition needed here from
$\mathrm{CovGmin.Hyp}$: on those rows
$\mathrm{laOf}(a,n)_{\tau i}=a(\tau,i)+n/2$, so its strict
inequality gives $a(\tau,i)>0$.
The original letter dictionary carries these degrees to the $p$-row
block on every component and for every $\mathrm{GoodEmb}$:
its differences $\mathrm{Dc.k}$ satisfy
$\mathrm{partOf}(\mathrm{Dc.k})=a$ after the specified relabelling.

An elementary consequence of the row-highest condition will be useful.
Let $X$ be the $p$-by-$n$ block of polynomial variables, keeping all
other variables fixed.  The original triangular action adds an earlier
row to a later row: $X_b\mapsto X_b+sX_a$ for $a<b$.
Its defining invariance for integral $s$ extends to every complex $s$
by polynomial identity.  The polynomial has positive degree in every
row.  If $X$ has rank less than $p$, choose a nonzero row relation and
its last nonzero coefficient.  Adding the earlier rows to that row
makes it zero by allowed triangular operations.  Homogeneity in that row
then gives $P(X)=0$.  In particular, for every nonzero $v\in\mathbb C^p$,
$P$ vanishes on matrices whose columns lie in $v^\perp$ for the bilinear
coordinate pairing.  Therefore $P(\partial)$ kills any holomorphic
function invariant under translation in the direction $v$ in each
column.  To see the last assertion directly, choose coordinates with
$v$ the last basis vector.  Restriction of the symbol $P$ to the
remaining derivative variables is zero, so each term contains a
derivative in one of the translation directions.

We check this for the \emph{actual} zero Fourier--Jacobi coefficient.
Fix any nonzero rational isotropic vector $\xi$.  The lattice
$E\xi\cap L$ is saturated and isotropic for the original unimodular
alternating form; successive symplectic completion gives an adapted
integral symplectic basis $(a_j,b_j)$.  The ordinary theta transformation
in this basis changes our fixed finite Schwartz vector into a finite
sum of rational characteristics, with its actual scalar factors.
For one such characteristic $(\alpha,\beta)$ put
\[
 \ell_{\tau j}=\tau\bigl(h(b_j,\xi)\bigr),\qquad
 C_{\tau,jk}=c_\tau^2\operatorname{Re}
                  (\ell_{\tau j}\overline{\ell_{\tau k}}),\qquad
 \nu_{m,\tau}=\frac{c_\tau^2}{4}
       \sum_{b=1}^n\left|\sum_j(m_{bj}+\alpha_{bj})\ell_{\tau j}\right|^2,
                                                        \tag{RT3}
\]
where $c_\tau=2|\tau(\xi_0)|$ is the original source constant.
The adapted period identities give
$\tau(x\cdot s)=\tau(x)+\sum_\tau s_\tau C_\tau$;
the matrices $h(x)$ and $Q(x)=\mathrm{splitMat}(x)$ are unchanged.
Thus the exact harmonic theta reading along the original flow is
\[
 \sum_m A_m\exp\left(4\pi i\sum_\tau\nu_{m,\tau}s_\tau\right),
 \qquad
 A_m=\det h(x)^{n/2}P(\partial_{\mathrm{letter}})F_m(0),
                                                        \tag{RT4}
\]
where
\[
 F_m(z)=\prod_{b=1}^n
       e^{\pi z_b^tQ(x)z_b/2}
       \theta_{m_b}\bigl(\tau(x),\alpha_b,\beta_b;h(x)z_b\bigr).
\]
Here $\theta_{m_b}$ is the single summand of the ordinary theta
series.  Its Gaussian majorant, also after each of the finitely many
derivatives in $P$, proves $\sum_m|A_m|<\infty$.
Harmonicity removes precisely the paired contractions in the original
split gauge, which is why \textup{(RT4)} is the reading of our section.

For completeness the permissible Borel gauge can always be chosen so
that $\mathrm{coper}(x,\xi)=0$.  In the two blocks write the period
vectors as $w,l$, with $\sum_iw_il_i=0$.
Their relevant projections are nonzero, since a nonzero isotropic
vector cannot lie in either definite summand.  Choose $d,e$ with
$\sum_{c\ge q}d_cw_c=1$ and $\sum_{a<q}e_al_a=1$, and put
$t=\sum_{a<q}w_al_a$.
The rectangular block
$K_{ac}=w_ad_c-e_al_c-te_ad_c$ satisfies
$\sum_{c\ge q}K_{ac}w_c=w_a$ and
$\sum_{a<q}K_{ac}l_a=-l_c$.
The corresponding upper unipotent frame change gives the asserted
gauge and has weight character one.
In this gauge the original $p$-row direction $v_\tau\ne0$, expressed
in the full period coordinates as $d_\tau$, satisfies
\[
       Q(x)d_\tau=0,\qquad
       (h(x)d_\tau)_j=-\frac{i c_\tau}{2}
                                      \overline{\ell_{\tau j}}.
                                                        \tag{RT5}
\]
These are the adapted period and adjoint identities; their signs and
the factor $1/2$ agree with the prescribed letter directions.

Here is a coordinate derivation of the identities used in
\textup{(RT3)--(RT5)}.  Write $\mathcal P_x$ and $\mathcal C_x$ for
the original period and co-period maps, including the $c_\tau$ factors
in $\mathcal C_x$.  Set
$A=(\mathcal P_x(a_j))_j$ and $B=(\mathcal P_x(b_j))_j$.
By definition
$h=A^{-1}$, $\tau=hB$ and
$Q=(\mathcal C_x(a_j))_jA^{-1}$.
The bilinear Riemann identity on the isotropic $a$-span makes $Q$
symmetric.  Let $d_\tau,e_\tau$ be respectively the $p$- and
$q$-block parts of $\mathcal P_x(\xi)$.  The eigenperiod convention is
\[
 \mathcal P_x(\alpha\xi)
       =\sum_\tau\bigl(\tau(\alpha)d_\tau
                             +\overline{\tau(\alpha)}e_\tau\bigr).
\]
The actual symplectic sign is $B_0(a_i,b_j)=-\delta_{ij}$.
Since $\alpha\xi$ is in the $a$-span, its $j$th $a$-coordinate
is $-B_0(\alpha\xi,b_j)$.  The original trace form gives
\[
 (h\mathcal P_x(\alpha\xi))_j
   =\sum_\tau\left(
      -\frac{i c_\tau}{2}\tau(\alpha)\overline{\ell_{\tau j}}
      +\frac{i c_\tau}{2}\overline{\tau(\alpha)}\ell_{\tau j}
                \right).                              \tag{RT7}
\]
Linear independence of the distinct complex embeddings on
$\mathcal O_E$ proves the two individual identities for $hd_\tau$
and $he_\tau$.  In the chosen gauge
$\mathcal C_x(\alpha\xi)=0$ for all $\alpha$, by its conjugate
eigencharacter rule.  The definition of $Q$ on the $a$-span then gives
$Q\mathcal P_x(\alpha\xi)=0$; the same independence gives
$Qd_\tau=Qe_\tau=0$.

Finally the original rank-one flow and its inverse are
$v\mapsto v-i c_\tau s_\tau h_\tau(v,\xi)\xi$ and the same
expression with $s_\tau$ negated.  Substitution in the two period
blocks gives
\[
 \mathcal P_{x\cdot s}(v)-\mathcal P_x(v)
   =\sum_\tau s_\tau\left(
       i c_\tau\tau(h(v,\xi))d_\tau
       -i c_\tau\overline{\tau(h(v,\xi))}e_\tau\right).
\]
The trace pairing and adaptation imply $h(a_j,\xi)=0$, so $A$ and
the co-period $a$-matrix, hence $h$ and $Q$, are fixed.
For $v=b_k$, applying $h$ and \textup{(RT7)} gives
$c_\tau^2(\ell_{\tau k}\overline{\ell_{\tau j}}+
\overline{\ell_{\tau k}}\ell_{\tau j})/2=C_{\tau,jk}$.
This proves the stated period shift.  For a real vector $u$,
$u^tC_\tau u=c_\tau^2|\sum_j u_j\ell_{\tau j}|^2$;
the ordinary theta exponent $\pi i u^t\tau u$ therefore gives
exactly the $4\pi i\nu_{m,\tau}$ in \textup{(RT4)}.

If $\nu_m=0$, nonnegativity in \textup{(RT3)} makes every displayed
linear form zero, as well as its complex conjugate.
Equations \textup{(RT5)} then make $F_m$ invariant under translation
by $d_\tau$ separately in each column.  The quadratic factor contributes
no residual linear or constant term because $Q(x)d_\tau=0$ and $Q$ is
symmetric.  The row argument therefore gives $A_m=0$.
The average of \textup{(RT4)} over expanding real boxes is zero:
each remaining frequency vector is nonzero and the absolutely summable
majorant permits passage to the limit.  Since the section is periodic
under the original arithmetic transvections, this average is its actual
zero Fourier--Jacobi coefficient.  This argument does not require each
individual characteristic summand to have the same period lattice.
It applies to every rational isotropic line, every component frame and
every original good embedding.\cite{raw-theta-cusp-source}

We also verify boundedness, which is a separate clause of the source
cusp space.  In an orthonormal frame let $\lambda$ be the shortest
nonzero period and $g=[E:\mathbb Q]r/2$.  The original probability
normalization gives squared $L^2$ norm one for the principal theta
section.  Projection of a Euclidean ball of fixed radius to the torus
has multiplicity at most $C_g(1+\lambda^{-1})^{2g}$,
by disjoint-ball packing in each translate of the period lattice.
The invariant metric in a fixed-radius holomorphic gauge is bounded
above and below by fixed Gaussian factors.  The holomorphic submean
inequality and Cauchy's inequality for the fixed finite jet of $P$
therefore give, for our finite sum of translated theta products,
\[
       \operatorname{pnorm}(\mathcal T(P),x)
           \le C(1+\lambda(x)^{-1})^{gn}.                \tag{RT6}
\]
All translations use the unitary theta metric; their finite number
changes $C$ only.  The covolume is fixed by the principal polarization,
so the passage between probability measure and Euclidean measure adds
only a fixed constant.  These period norms are the original positive
majorant norms, with their fixed $c_\tau$ factors retained.
In group coordinates $x=g_0x_0$,
$\lambda(x)\ge c/\|g_0^{-1}\|$ for the fixed period lattice.
Thus the associated automorphic function has moderate growth.

Here the rapid-decay theorem is applicable with all of its hypotheses.
The real group is the original
$\operatorname{Res}_{E^+/\mathbb Q}U(h)$, with compact real center
a product of norm-one circles.  Pass to the finite-index kernel of the
original finite theta character.  The function is smooth and
$K$-finite: on each compact full-flag fibre its right translates lie
in the finite-dimensional space of holomorphic sections of the fixed
weight line.  It is also $\mathfrak Z(\mathfrak g)$-finite.
Indeed the holomorphic Borel-equivariant frame reading is annihilated
by the right nilpotent Borel algebra and is an eigenvector for its
Cartan algebra; the Harish--Chandra projection of a central element
therefore acts by the corresponding weight scalar.
Finally every proper rational parabolic stabilizes an isotropic flag.
Choose $\xi$ in its first space $I_1$, of any dimension.
For the original purely imaginary parameter $t$, the transvection
$T_\xi(t)v=v+t h(v,\xi)\xi$ has
$(T_\xi(t)-1)V\subset I_1$ and fixes $I_1^\perp$.
It acts trivially on the associated graded flag, hence lies in the
unipotent radical.  That radical fixes $I_1$ pointwise, and
$uT_\xi(t)u^{-1}=T_{u\xi}(t)=T_\xi(t)$ proves centrality.
This includes $\dim I_1>1$.
Integration first over this rational central subgroup, whose average
is zero, proves the full parabolic constant term zero.
Miller--Schmid's Theorem~A now gives rapid decay
on every Siegel set.  Finitely many such sets cover the arithmetic
quotient, so the invariant norm is bounded everywhere.
This uses their central-finiteness hypothesis explicitly, as stated in
their equation~(6.1).\cite{raw-theta-rapid-decay}

If $p=0$ or $q=0$, the real unitary group and the full-flag quotient
are compact, and there is no nonzero rational isotropic vector.
Holomorphy gives boundedness directly, and the cusp condition is empty.
Thus the definite endpoints do not require the positive row argument.
In the indefinite range of $\mathrm{CovGmin.Hyp}$, the preceding
proof gives precisely membership in each original $X.\mathrm{SS}$.
Together with the algebraic section property and the intersection over
all $\mathrm{GoodEmb}$, the definition of $\mathrm{cuspW}$ gives
\[
       \mathcal T_{\Lambda,N,w}(P)
          \in\operatorname{cuspW}(\mathrm{wtZ}(a,n))
          =Y.V(\mathrm{wtZ}(a,n)).
\]
The constants in this cusp verification may depend on the fixed data,
level and polynomial.  No uniform archimedean norm estimate is asserted.

For an indicator $x_0$ the finite Fourier weight in this formula
is exactly $\zeta_M^{\langle x_0,q\rangle}$, with no additional
averaging factor.  Thus the constructed family meets the value
and arithmetic-lattice interfaces of the original raw generators.
Surjectivity onto the full holomorphic cusp space requires the
separate stable-range theta-spanning argument; neither this
construction nor the denominator estimate asserts that
surjectivity or the archimedean norm estimate.
